\documentclass[aspectratio=169,11pt]{beamer}
\usepackage[T1]{fontenc}
\usepackage{helvet}
\renewcommand{\familydefault}{\sfdefault}
\usepackage{booktabs,tikz,amsmath,amssymb,array}
\usepackage[most]{tcolorbox}

\definecolor{navy}{HTML}{1B2A41}
\definecolor{accent}{HTML}{C2410C}
\definecolor{tint}{HTML}{FBF3EA}
\definecolor{slate}{HTML}{F1F5F9}
\definecolor{good}{HTML}{15803D}
\definecolor{bad}{HTML}{B91C1C}
\definecolor{muted}{HTML}{6E798A}

\setbeamertemplate{navigation symbols}{}
\setbeamertemplate{footline}{}
\setbeamercolor{normal text}{fg=black}
\setbeamercolor{structure}{fg=navy}
\setbeamertemplate{itemize item}{\color{accent}$\blacktriangleright$}
\setbeamertemplate{itemize subitem}{\color{accent}\small$\triangleright$}
\setbeamertemplate{enumerate item}{\insertenumlabel.}
\setbeamersize{text margin left=36pt,text margin right=30pt}
\setbeamertemplate{frametitle}{%
  \nointerlineskip
  \begin{tikzpicture}[remember picture,overlay]
    \fill[navy] (current page.north west) rectangle ([yshift=-29.75pt]current page.north east);
    \node[anchor=west,text=white,font=\large\bfseries] at ([xshift=11.25pt,yshift=-14.9pt]current page.north west) {\insertframetitle};
  \end{tikzpicture}%
  \vspace*{22pt}}

\newtcolorbox{finding}{enhanced,colback=tint,colframe=tint,borderline west={2.2pt}{0pt}{accent},
  arc=0pt,outer arc=0pt,boxrule=0pt,left=7pt,right=6pt,top=4pt,bottom=4pt,fontupper=\small}
\newtcolorbox{idea}{enhanced,colback=slate,colframe=slate,borderline west={2.2pt}{0pt}{navy},
  arc=0pt,outer arc=0pt,boxrule=0pt,left=7pt,right=6pt,top=4pt,bottom=4pt,fontupper=\small}
\newcommand{\up}[1]{\textcolor{good}{#1}}
\newcommand{\dn}[1]{\textcolor{bad}{#1}}

\begin{document}

% ---------------------------------------------------------------- A (replaces p5)
\begin{frame}{Contributions}
\begin{enumerate}\setlength{\itemsep}{3pt}
\item A MeshGraphNet surrogate for the 1D viscous Burgers' equation (PDEBench).
\item A differentiable finite-difference residual for uniform and non-uniform meshes.
\item A decoder-only test-time adaptation procedure with anchor regularisation and best-state selection --- no labels used.
\item A \textbf{residual gate} that runs TTA only when the physics check can still see an error, calibrated on separate data.
\item A held-out evaluation: \textbf{20 unseen trajectories}, mean $\pm$ std, paired Wilcoxon tests; plus ablations and geometric shift.
\item A 2D extension on \textbf{trained} FNO-2D / U-Net-2D surrogates (8 datasets) and an interactive dashboard (Next.js + Express + PyTorch).
\end{enumerate}
\end{frame}

% ---------------------------------------------------------------- B (new, method)
\begin{frame}{Residual gate: adapt only when the physics sees an error}
\begin{columns}[T]
\begin{column}{0.52\textwidth}
\small
A low residual does not always mean a low error. If the model is already better than the residual can resolve, TTA only adds noise.
\vspace{4pt}
\begin{itemize}\setlength{\itemsep}{2pt}
\item Label-free score: \\[2pt]
  $\displaystyle q=\frac{R(\hat u)}{R(\text{copy } u_t)}$
\item Floor $q_{\text{floor}}$: median $q$ of the \emph{true} next state on calibration trajectories (per mesh).
\item Adapt iff $q > \kappa \cdot q_{\text{floor}}$.
\item $\kappa \in \{1,\dots,5,\infty\}$ chosen on calibration data; $\infty$ = never adapt.
\end{itemize}
\end{column}
\begin{column}{0.44\textwidth}
\begin{idea}
\textbf{Split (1D).} Train: trajectories 0--8999. Calibrate: 9900--9919. Test: 9980--9999. No test labels are used to pick $\kappa$.
\end{idea}
\vspace{2pt}
\begin{finding}
\textbf{Also fixed:} the best-state snapshot was saved \emph{after} the optimiser step it was scored on. It is now saved before, and the final state is checked too.
\end{finding}
\vspace{2pt}
{\footnotesize\color{muted} Crank--Nicolson residual: truth scores $0.18\times$ the copy residual ($0.77\times$ for implicit), so it is the sharper check.}
\end{column}
\end{columns}
\end{frame}

% ---------------------------------------------------------------- C (new, results headline)
\begin{frame}{Held-out evaluation: 20 unseen trajectories}
\vspace*{6pt}\begin{minipage}{\linewidth}\small 100 samples per mesh (20 trajectories $\times$ 5 timesteps). Original model, 12 steps, $\eta=10^{-4}$, $\lambda=0.1$.\end{minipage}
\vspace{4pt}
\begin{center}\resizebox{\linewidth}{!}{%
\begin{tabular}{l c c c c c}
\toprule
Mesh & Baseline rel.\ $L_2$ & Plain TTA & \textbf{Gated TTA} & Worse (plain$\to$gated) & $p$ (gated) \\
\midrule
Uniform & 9.62\% & \up{+15.9\%} $\pm$ 19.4 & \textbf{\up{+18.9\%}} $\pm$ 15.4 & 11\% $\to$ 4\% & $1.3\times10^{-5}$ \\
Jitter $\pm35\%$ & 9.60\% & \up{+14.4\%} $\pm$ 25.4 & \textbf{\up{+18.1\%}} $\pm$ 13.5 & 8\% $\to$ 0\% & $2\times10^{-6}$ \\
Warped & 9.49\% & \up{+14.0\%} $\pm$ 26.9 & \textbf{\up{+18.0\%}} $\pm$ 16.1 & 8\% $\to$ 0\% & $2\times10^{-6}$ \\
Coarse $N{=}512$ & 9.57\% & \up{+22.7\%} $\pm$ 16.1 & \textbf{\up{+22.7\%}} $\pm$ 16.1 & 2\% $\to$ 2\% & $2\times10^{-6}$ \\
\bottomrule
\end{tabular}}
\end{center}
\vspace{2pt}
\begin{finding}
Every gated gain is significant (paired two-sided Wilcoxon over trajectories, $p\le1.3\times10^{-5}$), and the gate also \textbf{cuts the spread}: std 13--16\% vs.\ 19--27\% for plain TTA. Earlier we quoted 73.9\% from 3 hand-picked samples; \emph{these} are the numbers to quote.
\end{finding}
{\footnotesize\color{muted} Error reduction = MSE change per trajectory, mean $\pm$ std over the 20 trajectories (positive = TTA helped).}
\end{frame}

% ---------------------------------------------------------------- D (replaces p20)
\begin{frame}{Geometric shift and an accurate model}
\begin{columns}[T]
\begin{column}{0.47\textwidth}
\small
\textbf{Geometric shift is no longer an open problem.} The earlier losses ($+13.8\%$ MSE on warped) came from the best-state bug. After the fix, plain TTA on the original model:
\vspace{3pt}
\begin{center}
\begin{tabular}{lc}
\toprule
Shift & Error reduction \\
\midrule
Jitter $\pm35\%$ & \up{+34.9\%} \\
Warped & \up{+28.0\%} \\
Coarse ($N{=}512$) & \up{+36.9\%} \\
\bottomrule
\end{tabular}
\end{center}
{\footnotesize\color{muted} Same samples as the original table; the held-out table on slide 10 is the reliable version.}
\end{column}
\begin{column}{0.49\textwidth}
\small
\textbf{v2 model} (delta prediction, normalisation; 2.35\% one-step error vs.\ 12.8\%):
\vspace{3pt}
\begin{center}\footnotesize
\begin{tabular}{lcc}
\toprule
Mesh & Plain TTA & \textbf{Gated} \\
\midrule
Uniform & \dn{$-190.6\%$} & $-8.2\%$ {\scriptsize($p{=}0.06$)} \\
Jitter & \dn{$-61.9\%$} & 0\% {\scriptsize(shut)} \\
Warped & \up{$+13.7\%$} & \textbf{\up{+13.1\%}} \\
Coarse & \up{$+36.9\%$} & \textbf{\up{+36.9\%}} \\
\bottomrule
\end{tabular}
\end{center}
\begin{finding}
Plain TTA hurts 82--89\% of v2 samples on uniform and jitter; the gate cuts that to $\le$5\% and keeps the gains where v2 is weak.
\end{finding}
\end{column}
\end{columns}
\end{frame}

% ---------------------------------------------------------------- E (replaces p22)
\begin{frame}{2D extension: flow past obstacles}
\begin{columns}[T]
\begin{column}{0.49\textwidth}
\textbf{Data --- Lattice Boltzmann (D2Q9, BGK)}
\begin{itemize}\small\setlength{\itemsep}{1pt}
\item 9 particle populations per cell; collide $f\leftarrow f-\omega(f-f^{eq})$, then stream
\item Bounce-back walls on obstacle masks; $Re\approx 60$--$195$
\item $48\times96$ grid: cylinder, NACA\,0012, square, backward step, twin cylinders
\item Plus heat, advection, NS vorticity ($64\times64$)
\end{itemize}
\textbf{Surrogates}
\begin{itemize}\small\setlength{\itemsep}{1pt}
\item FNO-2D (spectral, 1.19\,M params), U-Net-2D (0.47\,M)
\end{itemize}
\end{column}
\begin{column}{0.47\textwidth}
\textbf{2D residuals} (conv-kernel stencils)
{\small\[\begin{gathered}
\nabla\cdot\mathbf u=\partial_x u_x+\partial_y u_y=0\\
\partial_t u-\kappa\nabla^2u=0\\
\partial_t\mathbf u+(\mathbf u\cdot\nabla)\mathbf u+\nabla p-\nu\nabla^2\mathbf u=0
\end{gathered}\]}
\begin{idea}
\textbf{Trained and benchmarked.} 10 trajectories per dataset: 6 train, 1 calibration (traj.\ 9), 3 test (2, 5, 8). On the LBM flows the models are \textbf{2--13$\times$ more accurate} than copying the previous frame (e.g.\ square: 0.08\% vs.\ 0.78\%).
\end{idea}
\end{column}
\end{columns}
\end{frame}

% ---------------------------------------------------------------- F (new, 2D results)
\begin{frame}{2D results: the gate decides where TTA is safe}
\vspace*{6pt}\begin{minipage}{\linewidth}\small Trained FNO-2D / U-Net-2D on 8 datasets (5 Lattice-Boltzmann obstacle flows, heat, advection, vorticity). Per dataset: 6 train, 1 calibration, 3 test trajectories (57 test frames); $\kappa$ and $q_{\text{floor}}$ from the calibration trajectory only.\end{minipage}
\vspace{3pt}
\begin{center}\resizebox{\linewidth}{!}{%
\begin{tabular}{l l c c c}
\toprule
Dataset & Model & Plain TTA & \textbf{Gated TTA} & Frames worse (plain$\to$gated) \\
\midrule
5 LBM obstacle flows & FNO-2D & $+0.9\%$ to \dn{$-20.2\%$} & 0\% (gate shut) & 35--100\% $\to$ \textbf{0\%} \\
5 LBM obstacle flows & U-Net-2D & \dn{$-11.9\%$ to $-249\%$} & 0\% (gate shut) & 55--100\% $\to$ \textbf{0\%} \\
\midrule
Heat & FNO / U-Net & \up{+25.0\% / +45.4\%} & \textbf{\up{+25.0\% / +45.4\%}} & 0\% $\to$ 0\% \\
Advection & FNO / U-Net & \up{+6.6\% / +67.5\%} & \textbf{\up{+6.6\% / +67.5\%}} & 0\% $\to$ 0\% \\
Vorticity & FNO / U-Net & \up{+5.7\% / +57.7\%} & \textbf{\up{+5.7\% / +57.7\%}} & 0\% $\to$ 0\% \\
\bottomrule
\end{tabular}}
\end{center}
\vspace{2pt}
\begin{finding}
The gate opened on \textbf{every} frame of the three equations with an exact residual ($p\approx5\times10^{-11}$ over frames) and on \textbf{no} frame of the LBM flows, where the trained models already beat what the residual can resolve. No 2D test frame got worse.
\end{finding}
{\footnotesize\color{muted} Frames within a trajectory are correlated, so treat 2D $p$-values as indicative.}
\end{frame}

% ---------------------------------------------------------------- G (replaces p25)
\begin{frame}{Limitations \& future work}
\begin{columns}[T]
\begin{column}{0.48\textwidth}
\textbf{Limitations}
\begin{itemize}\small\setlength{\itemsep}{2pt}
\item Gains depend on the residual: TTA only helps where the PDE check is sharper than the model's error.
\item The gate protects, but it cannot make TTA help an already-accurate model (v2 uniform, LBM flows).
\item Needs the governing PDE and its parameters.
\item $\sim$0.35\,s overhead per sample on CPU.
\item 2D tests: 3 trajectories; frames correlated.
\item Single-step evaluation; v2 rollouts unstable.
\end{itemize}
\end{column}
\begin{column}{0.48\textwidth}
\textbf{Future work}
\begin{itemize}\small\setlength{\itemsep}{2pt}
\item Sharper residuals for LBM flows (e.g.\ lattice-level consistency) so the gate can open there.
\item More 2D test trajectories for trajectory-level tests.
\item Shift in the data (true $\nu=0.002$ simulations).
\item Trajectory-level TTA for long autoregressive rollouts.
\item Adapt edge encoder / LayerNorm under geometric shift.
\end{itemize}
\end{column}
\end{columns}
\end{frame}

% ---------------------------------------------------------------- H (replaces p26)
\begin{frame}{Conclusion}
\begin{center}
\textbf{A trained surrogate can be corrected at test time without labels by minimising the PDE residual through its decoder --- and a label-free gate tells us when not to.}
\end{center}
\vspace{4pt}
\begin{itemize}\small\setlength{\itemsep}{3pt}
\item \textbf{Held-out 1D:} +18--23\% lower error on all four meshes (20 unseen trajectories, $p\le1.3\times10^{-5}$), with $\le$4\% of samples made worse.
\item \textbf{Geometric shift:} works after the best-state fix (+28--37\%).
\item \textbf{2D:} +45--68\% (U-Net) on heat, advection and vorticity; the gate stays shut on obstacle flows, so no frame got worse.
\item Decoder-only adaptation is 4$\times$ cheaper and far more stable than full fine-tuning.
\item TTA is only as good as its physics residual --- the gate measures that directly.
\end{itemize}
\end{frame}

\end{document}
